\IfFileExists{stacks-project.cls}{%
\documentclass{stacks-project}
}{%
\documentclass{amsart}
}

% For dealing with references we use the comment environment
\usepackage{verbatim}
\newenvironment{reference}{\comment}{\endcomment}
%\newenvironment{reference}{}{}
\newenvironment{slogan}{\comment}{\endcomment}
\newenvironment{history}{\comment}{\endcomment}

% For commutative diagrams we use Xy-pic
\usepackage[all]{xy}

% We use 2cell for 2-commutative diagrams.
\xyoption{2cell}
\UseAllTwocells

% We use multicol for the list of chapters between chapters
\usepackage{multicol}

% This is generally recommended for better output
\usepackage{lmodern}
\usepackage[T1]{fontenc}

% For cross-file-references
\usepackage{xr-hyper}

% Package for hypertext links:
\usepackage{hyperref}

% For any local file, say "hello.tex" you want to link to please
% use \externaldocument[hello-]{hello}
\externaldocument[introduction-]{introduction}
\externaldocument[conventions-]{conventions}
\externaldocument[sets-]{sets}
\externaldocument[categories-]{categories}
\externaldocument[topology-]{topology}
\externaldocument[sheaves-]{sheaves}
\externaldocument[sites-]{sites}
\externaldocument[stacks-]{stacks}
\externaldocument[fields-]{fields}
\externaldocument[algebra-]{algebra}
\externaldocument[brauer-]{brauer}
\externaldocument[homology-]{homology}
\externaldocument[derived-]{derived}
\externaldocument[simplicial-]{simplicial}
\externaldocument[more-algebra-]{more-algebra}
\externaldocument[smoothing-]{smoothing}
\externaldocument[modules-]{modules}
\externaldocument[sites-modules-]{sites-modules}
\externaldocument[injectives-]{injectives}
\externaldocument[cohomology-]{cohomology}
\externaldocument[sites-cohomology-]{sites-cohomology}
\externaldocument[dga-]{dga}
\externaldocument[dpa-]{dpa}
\externaldocument[sdga-]{sdga}
\externaldocument[hypercovering-]{hypercovering}
\externaldocument[schemes-]{schemes}
\externaldocument[constructions-]{constructions}
\externaldocument[properties-]{properties}
\externaldocument[morphisms-]{morphisms}
\externaldocument[coherent-]{coherent}
\externaldocument[divisors-]{divisors}
\externaldocument[limits-]{limits}
\externaldocument[varieties-]{varieties}
\externaldocument[topologies-]{topologies}
\externaldocument[descent-]{descent}
\externaldocument[perfect-]{perfect}
\externaldocument[more-morphisms-]{more-morphisms}
\externaldocument[flat-]{flat}
\externaldocument[groupoids-]{groupoids}
\externaldocument[more-groupoids-]{more-groupoids}
\externaldocument[etale-]{etale}
\externaldocument[chow-]{chow}
\externaldocument[intersection-]{intersection}
\externaldocument[pic-]{pic}
\externaldocument[weil-]{weil}
\externaldocument[adequate-]{adequate}
\externaldocument[dualizing-]{dualizing}
\externaldocument[duality-]{duality}
\externaldocument[discriminant-]{discriminant}
\externaldocument[derham-]{derham}
\externaldocument[local-cohomology-]{local-cohomology}
\externaldocument[algebraization-]{algebraization}
\externaldocument[curves-]{curves}
\externaldocument[resolve-]{resolve}
\externaldocument[models-]{models}
\externaldocument[functors-]{functors}
\externaldocument[equiv-]{equiv}
\externaldocument[pione-]{pione}
\externaldocument[etale-cohomology-]{etale-cohomology}
\externaldocument[proetale-]{proetale}
\externaldocument[relative-cycles-]{relative-cycles}
\externaldocument[more-etale-]{more-etale}
\externaldocument[trace-]{trace}
\externaldocument[crystalline-]{crystalline}
\externaldocument[spaces-]{spaces}
\externaldocument[spaces-properties-]{spaces-properties}
\externaldocument[spaces-morphisms-]{spaces-morphisms}
\externaldocument[decent-spaces-]{decent-spaces}
\externaldocument[spaces-cohomology-]{spaces-cohomology}
\externaldocument[spaces-limits-]{spaces-limits}
\externaldocument[spaces-divisors-]{spaces-divisors}
\externaldocument[spaces-over-fields-]{spaces-over-fields}
\externaldocument[spaces-topologies-]{spaces-topologies}
\externaldocument[spaces-descent-]{spaces-descent}
\externaldocument[spaces-perfect-]{spaces-perfect}
\externaldocument[spaces-more-morphisms-]{spaces-more-morphisms}
\externaldocument[spaces-flat-]{spaces-flat}
\externaldocument[spaces-groupoids-]{spaces-groupoids}
\externaldocument[spaces-more-groupoids-]{spaces-more-groupoids}
\externaldocument[bootstrap-]{bootstrap}
\externaldocument[spaces-pushouts-]{spaces-pushouts}
\externaldocument[spaces-chow-]{spaces-chow}
\externaldocument[groupoids-quotients-]{groupoids-quotients}
\externaldocument[spaces-more-cohomology-]{spaces-more-cohomology}
\externaldocument[spaces-simplicial-]{spaces-simplicial}
\externaldocument[spaces-duality-]{spaces-duality}
\externaldocument[formal-spaces-]{formal-spaces}
\externaldocument[restricted-]{restricted}
\externaldocument[spaces-resolve-]{spaces-resolve}
\externaldocument[formal-defos-]{formal-defos}
\externaldocument[defos-]{defos}
\externaldocument[cotangent-]{cotangent}
\externaldocument[examples-defos-]{examples-defos}
\externaldocument[algebraic-]{algebraic}
\externaldocument[examples-stacks-]{examples-stacks}
\externaldocument[stacks-sheaves-]{stacks-sheaves}
\externaldocument[criteria-]{criteria}
\externaldocument[artin-]{artin}
\externaldocument[quot-]{quot}
\externaldocument[stacks-properties-]{stacks-properties}
\externaldocument[stacks-morphisms-]{stacks-morphisms}
\externaldocument[stacks-limits-]{stacks-limits}
\externaldocument[stacks-cohomology-]{stacks-cohomology}
\externaldocument[stacks-perfect-]{stacks-perfect}
\externaldocument[stacks-introduction-]{stacks-introduction}
\externaldocument[stacks-more-morphisms-]{stacks-more-morphisms}
\externaldocument[stacks-geometry-]{stacks-geometry}
\externaldocument[moduli-]{moduli}
\externaldocument[moduli-curves-]{moduli-curves}
\externaldocument[examples-]{examples}
\externaldocument[exercises-]{exercises}
\externaldocument[guide-]{guide}
\externaldocument[desirables-]{desirables}
\externaldocument[coding-]{coding}
\externaldocument[obsolete-]{obsolete}
\externaldocument[fdl-]{fdl}
\externaldocument[index-]{index}

% Theorem environments.
%
\theoremstyle{plain}
\newtheorem{theorem}[subsection]{Theorem}
\newtheorem{proposition}[subsection]{Proposition}
\newtheorem{lemma}[subsection]{Lemma}

\theoremstyle{definition}
\newtheorem{definition}[subsection]{Definition}
\newtheorem{example}[subsection]{Example}
\newtheorem{exercise}[subsection]{Exercise}
\newtheorem{situation}[subsection]{Situation}

\theoremstyle{remark}
\newtheorem{remark}[subsection]{Remark}
\newtheorem{remarks}[subsection]{Remarks}

\numberwithin{equation}{subsection}

% Macros
%
\def\lim{\mathop{\mathrm{lim}}\nolimits}
\def\colim{\mathop{\mathrm{colim}}\nolimits}
\def\Spec{\mathop{\mathrm{Spec}}}
\def\Hom{\mathop{\mathrm{Hom}}\nolimits}
\def\Ext{\mathop{\mathrm{Ext}}\nolimits}
\def\SheafHom{\mathop{\mathcal{H}\!\mathit{om}}\nolimits}
\def\SheafExt{\mathop{\mathcal{E}\!\mathit{xt}}\nolimits}
\def\Sch{\mathit{Sch}}
\def\Mor{\mathop{\mathrm{Mor}}\nolimits}
\def\Ob{\mathop{\mathrm{Ob}}\nolimits}
\def\Sh{\mathop{\mathit{Sh}}\nolimits}
\def\NL{\mathop{N\!L}\nolimits}
\def\CH{\mathop{\mathrm{CH}}\nolimits}
\def\proetale{{pro\text{-}\acute{e}tale}}
\def\etale{{\acute{e}tale}}
\def\QCoh{\mathit{QCoh}}
\def\Ker{\mathop{\mathrm{Ker}}}
\def\Im{\mathop{\mathrm{Im}}}
\def\Coker{\mathop{\mathrm{Coker}}}
\def\Coim{\mathop{\mathrm{Coim}}}

% Boxtimes
%
\DeclareMathSymbol{\boxtimes}{\mathbin}{AMSa}{"02}

%
% Macros for moduli stacks/spaces
%
\def\QCohstack{\mathcal{QC}\!\mathit{oh}}
\def\Cohstack{\mathcal{C}\!\mathit{oh}}
\def\Spacesstack{\mathcal{S}\!\mathit{paces}}
\def\Quotfunctor{\mathrm{Quot}}
\def\Hilbfunctor{\mathrm{Hilb}}
\def\Curvesstack{\mathcal{C}\!\mathit{urves}}
\def\Polarizedstack{\mathcal{P}\!\mathit{olarized}}
\def\Complexesstack{\mathcal{C}\!\mathit{omplexes}}
% \Pic is the operator that assigns to X its picard group, usage \Pic(X)
% \Picardstack_{X/B} denotes the Picard stack of X over B
% \Picardfunctor_{X/B} denotes the Picard functor of X over B
\def\Pic{\mathop{\mathrm{Pic}}\nolimits}
\def\Picardstack{\mathcal{P}\!\mathit{ic}}
\def\Picardfunctor{\mathrm{Pic}}
\def\Deformationcategory{\mathcal{D}\!\mathit{ef}}
\usepackage{cleveref}
\crefname{originallemma}{original lemma}{original lemmas}
\Crefname{originallemma}{Original lemma}{Original lemmas}
\makeatletter
\newlabel{derived-lemma-find-existence-computes}{{06XN}{source}{}{sourcepointer-1}{}}
\newlabel{derived-lemma-find-existence-computes@cref}{{[originallemma][0][]06XN}{[1][0][]source}{}{}{}}
\makeatother
\crefname{originallemma}{original lemma}{original lemmas}
\Crefname{originallemma}{Original lemma}{Original lemmas}
\makeatletter
\newlabel{homology-lemma-good-resolution-gives-qis}{{09IZ}{source}{}{sourcepointer-2}{}}
\newlabel{homology-lemma-good-resolution-gives-qis@cref}{{[originallemma][0][]09IZ}{[1][0][]source}{}{}{}}
\makeatother
\crefname{originallemma}{original lemma}{original lemmas}
\Crefname{originallemma}{Original lemma}{Original lemmas}
\makeatletter
\newlabel{derived-lemma-representable-homological}{{0149}{source}{}{sourcepointer-3}{}}
\newlabel{derived-lemma-representable-homological@cref}{{[originallemma][0][]0149}{[1][0][]source}{}{}{}}
\makeatother
\crefname{originallemma}{original lemma}{original lemmas}
\Crefname{originallemma}{Original lemma}{Original lemmas}
\makeatletter
\newlabel{homology-lemma-good-right-resolution-gives-qis}{{09J0}{source}{}{sourcepointer-4}{}}
\newlabel{homology-lemma-good-right-resolution-gives-qis@cref}{{[originallemma][0][]09J0}{[1][0][]source}{}{}{}}
\makeatother
\crefname{originallemma}{original lemma}{original lemmas}
\Crefname{originallemma}{Original lemma}{Original lemmas}
\makeatletter
\newlabel{derived-lemma-easier-axiom-four}{{014C}{source}{}{sourcepointer-5}{}}
\newlabel{derived-lemma-easier-axiom-four@cref}{{[originallemma][0][]014C}{[1][0][]source}{}{}{}}
\makeatother
\crefname{originallemma}{original lemma}{original lemmas}
\Crefname{originallemma}{Original lemma}{Original lemmas}
\makeatletter
\newlabel{derived-lemma-make-injective}{{013N}{source}{}{sourcepointer-6}{}}
\newlabel{derived-lemma-make-injective@cref}{{[originallemma][0][]013N}{[1][0][]source}{}{}{}}
\makeatother
\title{The derived tensor functor with a compact complex is fully faithful}
\author{The Stacks Project Authors}
\date{stacks-project:dga.tex:local-snapshot-20261012; GNU FDL 1.2 or later}
\begin{document}
\maketitle
\noindent\textit{Editorial note.} The counted statement is the author's own lemma environment labelled lemma-fully-faithful-in-compact-case (source lines 6338-6356); the proof environment begins on the second line after it and is its own proof, closing at source line 6374.\par\medskip
\noindent\textit{Editorial note.} The statement and the proof invoke three author results of the same chapter by label; all three are retained as context spans and the generator expands the context transitively (budget 800 lines for this card), while reference targets in other chapters that carry canonical project tags are auto-declared as non-load-bearing source pointers.\par\medskip
\noindent\textit{Editorial note.} The retained preamble is the chapter's own minimal preamble (sources/stacks-dga/preamble.tex, taken with its real line count), a verbatim excerpt of the native Xy-pic setup, the native theorem declarations with the mathematical macros, and the project's own no-op reference environments; no macro was invented or redefined.\par\medskip
\noindent\textit{Editorial scope.} One original result of the Stacks Project's Differential Graded Algebra chapter, counted once: the derived tensor functor with a compact complex whose endomorphism algebra computes its Ext groups is fully faithful.\par\medskip
\noindent\textit{Allowed standard prerequisites (editorial).}\begin{itemize}
\item Ringed site and complexes of O-modules. A ringed site (C, O) and the derived category D(O) of complexes of O-modules, with its compact objects.. It is the setting and the hypothesis that K\textasciicircum{}bullet represents a compact object (source lines 6340-6344).
\item Derived tensor product with a fixed complex. The left derived functor - tensor\textasciicircum{}L\_E K\textasciicircum{}bullet : D(E,d) to D(O) of the earlier lemma, and its right adjoint R Hom(K\textasciicircum{}bullet, -) (the author's lemmas lemma-tensor-with-complex-derived and lemma-tensor-with-complex-hom-adjoint, retained as context).. The functor whose full faithfulness is asserted, and the adjunction used to reduce the claim (source lines 6350-6360).
\item Endomorphism dg algebra computing Ext. E = Hom over the dg category of complexes from K\textasciicircum{}bullet to itself, assumed to compute the Ext groups of K\textasciicircum{}bullet in D(O).. It is the algebra the source category D(E,d) is taken over, and it is the second hypothesis of the statement (source lines 6345-6346, 6372-6373).
\item Property (P) modules and their filtration sequences. Every dg E-module has a quasi-isomorphic representative with property (P), carrying an admissible filtration with graded pieces direct sums of shifts of E, and a split admissible short exact sequence displaying the filtration (the author's lemma lemma-property-P-sequence, retained as context).. The replacements and reductions of the proof run along this filtration (source lines 6366-6371).
\end{itemize}\medskip
\ifdefined\hypertarget\hypertarget{sourcepointer-1}{}\fi
\noindent\textit{Source reference.} Original lemma 06XN: sources/stacks-chow/tags.txt, source lines 8857--8857. Cross-chapter pointer whose target carries the canonical tag 06XN in the project tag list; declared so the reference resolves inside the assembled package, while every step the counted argument uses remains inside the retained spans.\par\medskip
\ifdefined\hypertarget\hypertarget{sourcepointer-2}{}\fi
\noindent\textit{Source reference.} Original lemma 09IZ: sources/stacks-chow/tags.txt, source lines 12057--12057. Cross-chapter pointer whose target carries the canonical tag 09IZ in the project tag list; declared so the reference resolves inside the assembled package, while every step the counted argument uses remains inside the retained spans.\par\medskip
\ifdefined\hypertarget\hypertarget{sourcepointer-3}{}\fi
\noindent\textit{Source reference.} Original lemma 0149: sources/stacks-chow/tags.txt, source lines 1464--1464. Cross-chapter pointer whose target carries the canonical tag 0149 in the project tag list; declared so the reference resolves inside the assembled package, while every step the counted argument uses remains inside the retained spans.\par\medskip
\ifdefined\hypertarget\hypertarget{sourcepointer-4}{}\fi
\noindent\textit{Source reference.} Original lemma 09J0: sources/stacks-chow/tags.txt, source lines 12058--12058. Cross-chapter pointer whose target carries the canonical tag 09J0 in the project tag list; declared so the reference resolves inside the assembled package, while every step the counted argument uses remains inside the retained spans.\par\medskip
\ifdefined\hypertarget\hypertarget{sourcepointer-5}{}\fi
\noindent\textit{Source reference.} Original lemma 014C: sources/stacks-chow/tags.txt, source lines 1467--1467. Cross-chapter pointer whose target carries the canonical tag 014C in the project tag list; declared so the reference resolves inside the assembled package, while every step the counted argument uses remains inside the retained spans.\par\medskip
\ifdefined\hypertarget\hypertarget{sourcepointer-6}{}\fi
\noindent\textit{Source reference.} Original lemma 013N: sources/stacks-chow/tags.txt, source lines 1442--1442. Cross-chapter pointer whose target carries the canonical tag 013N in the project tag list; declared so the reference resolves inside the assembled package, while every step the counted argument uses remains inside the retained spans.\par\medskip
\begin{lemma}
\label{lemma-tensor-with-complex-derived}
The functor $K(\text{Mod}_{(E, \text{d})}) \to K(\mathcal{O})$
of Lemma \ref{lemma-tensor-with-complex-homotopy} has a left derived
version defined on all of $D(E, \text{d})$. We denote it
$- \otimes_E^\mathbf{L} K^\bullet : D(E, \text{d}) \to D(\mathcal{O})$.
\end{lemma}

\begin{proof}
We will use
Derived Categories, Lemma \ref{derived-lemma-find-existence-computes}
to prove this. As our collection $\mathcal{P}$
of objects we will use the objects with property (P).
Property (1) was shown in Lemma \ref{lemma-resolve}.
Property (2) holds because if $s : P \to P'$ is a quasi-isomorphism
of modules with property (P), then $s$ is a homotopy equivalence
by Lemma \ref{lemma-hom-derived}.
\end{proof}

\begin{lemma}
\label{lemma-tensor-with-complex-hom-adjoint}
Let $(\mathcal{C}, \mathcal{O})$ be a ringed site.
Let $K^\bullet$ be a complex of $\mathcal{O}$-modules.
Then the functor
$$
- \otimes_E^\mathbf{L} K^\bullet :
D(E, \text{d})
\longrightarrow
D(\mathcal{O})
$$
of Lemma \ref{lemma-tensor-with-complex-derived} is a left adjoint
of the functor
$$
R\Hom(K^\bullet, -) : D(\mathcal{O}) \longrightarrow D(E, \text{d})
$$
of Lemma \ref{lemma-existence-of-derived}.
\end{lemma}

\begin{proof}
The statement means that we have
$$
\Hom_{D(E, \text{d})}(M, R\Hom(K^\bullet, L^\bullet)) =
\Hom_{D(\mathcal{O})}(M \otimes^\mathbf{L}_E K^\bullet, L^\bullet)
$$
bifunctorially in $M$ and $L^\bullet$. To see this we may replace $M$
by a differential graded $E$-module $P$ with property (P).
We also may replace $L^\bullet$ by a K-injective complex of
$\mathcal{O}$-modules $I^\bullet$. The computation
of the derived functors given in the lemmas referenced in the statement
combined with Lemma \ref{lemma-hom-derived} translates the above into
$$
\Hom_{K(\text{Mod}_{(E, \text{d})})}
(P, \Hom_\mathcal{B}(K^\bullet, I^\bullet)) =
\Hom_{K(\mathcal{O})}(P \otimes_E K^\bullet, I^\bullet)
$$
where $\mathcal{B} = \text{Comp}^{dg}(\mathcal{O})$. 
There is an evaluation map from right to left functorial
in $P$ and $I^\bullet$ (details omitted).
Choose a filtration $F_\bullet$ on $P$ as in the definition of property (P).
By Lemma \ref{lemma-property-P-sequence} and the fact that
both sides of the equation are homological functors in $P$
on $K(\text{Mod}_{(E, \text{d})})$
we reduce to the case where $P$ is replaced by
the differential graded $E$-module $\bigoplus F_iP$.
Since both sides turn direct sums in the variable $P$
into direct products we reduce to the case where $P$ is one of the
differential graded $E$-modules $F_iP$.
Since each $F_iP$ has a finite filtration (given by admissible
monomorphisms) whose graded pieces are graded projective $E$-modules
we reduce to the case where $P$ is a graded projective $E$-module.
In this case we clearly have
$$
\Hom_{\text{Mod}^{dg}_{(E, \text{d})}}
(P, \Hom_\mathcal{B}(K^\bullet, I^\bullet)) =
\Hom_{\text{Comp}^{dg}(\mathcal{O})}(P \otimes_E K^\bullet, I^\bullet)
$$
as graded $\mathbf{Z}$-modules (because this statement reduces to the case
$P = E[k]$ where it is obvious). As the isomorphism is compatible with
differentials we conclude.
\end{proof}

\begin{lemma}
\label{lemma-property-P-sequence}
Let $(A, \text{d})$ be a differential graded algebra.
Let $P$ be a differential graded $A$-module. If $F_\bullet$
is a filtration as in property (P), then we obtain an
admissible short exact sequence
$$
0 \to
\bigoplus\nolimits F_iP \to
\bigoplus\nolimits F_iP \to P \to 0
$$
of differential graded $A$-modules.
\end{lemma}

\begin{proof}
The second map is the direct sum of the inclusion maps.
The first map on the summand $F_iP$ of the source is the sum
of the identity $F_iP \to F_iP$ and the negative of the inclusion
map $F_iP \to F_{i + 1}P$. Choose homomorphisms $s_i : F_{i + 1}P \to F_iP$
of graded $A$-modules which are left inverse to the inclusion maps.
Composing gives maps $s_{j, i} : F_jP \to F_iP$ for all $j > i$.
Then a left inverse of the first arrow maps $x \in F_jP$ to
$(s_{j, 0}(x), s_{j, 1}(x), \ldots, s_{j, j - 1}(x), 0, \ldots)$
in $\bigoplus F_iP$.
\end{proof}

\begin{lemma}
\label{lemma-hom-derived}
Let $(A, \text{d})$ be a differential graded algebra.
Let $M$ and $N$ be differential graded $A$-modules.
\begin{enumerate}
\item Let $P \to M$ be a P-resolution as in
Lemma \ref{lemma-resolve}. Then
$$
\Hom_{D(A, \text{d})}(M, N) =
\Hom_{K(\text{Mod}_{(A, \text{d})})}(P, N)
$$
\item Let $N \to I$ be an I-resolution as in
Lemma \ref{lemma-right-resolution}. Then
$$
\Hom_{D(A, \text{d})}(M, N) =
\Hom_{K(\text{Mod}_{(A, \text{d})})}(M, I)
$$
\end{enumerate}
\end{lemma}

\begin{proof}
Let $P \to M$ be as in (1). Since $P \to M$ is a quasi-isomorphism we see that
$$
\Hom_{D(A, \text{d})}(P, N) = \Hom_{D(A, \text{d})}(M, N)
$$
by definition of the derived category. A morphism
$f : P \to N$ in $D(A, \text{d})$ is equal to
$s^{-1}f'$ where $f' : P \to N'$ is a morphism and
$s : N \to N'$ is a quasi-isomorphism. Choose a distinguished triangle
$$
N \to N' \to Q \to N[1]
$$
As $s$ is a quasi-isomorphism, we see that $Q$ is acyclic. Thus
$\Hom_{K(\text{Mod}_{(A, \text{d})})}(P, Q[k]) = 0$ for all $k$ by
Lemma \ref{lemma-property-P-K-projective}. Since
$\Hom_{K(\text{Mod}_{(A, \text{d})})}(P, -)$
is cohomological, we conclude that we can lift $f' : P \to N'$
uniquely to a morphism $f : P \to N$. This finishes the proof.

\medskip\noindent
The proof of (2) is dual to that of (1) using
Lemma \ref{lemma-property-I-K-injective} in stead of
Lemma \ref{lemma-property-P-K-projective}.
\end{proof}

\begin{lemma}
\label{lemma-resolve}
Let $(A, \text{d})$ be a differential graded algebra.
Let $M$ be a differential graded $A$-module. There exists a homomorphism
$P \to M$ of differential graded $A$-modules such that
\begin{enumerate}
\item $P \to M$ is a quasi-isomorphism, and
\item $P$ has property (P).
\end{enumerate}
\end{lemma}

\begin{proof}
Set $M = M_0$. We inductively choose short exact sequences
$$
0 \to M_{i + 1} \to P_i \to M_i \to 0
$$
where the maps $P_i \to M_i$ are chosen as in Lemma \ref{lemma-good-quotient}.
This gives a ``resolution''
$$
\ldots \to P_2 \xrightarrow{f_2} P_1 \xrightarrow{f_1} P_0 \to M \to 0
$$
Then we set
$$
P = \bigoplus\nolimits_{i \geq 0} P_i
$$
as an $A$-module with grading given by
$P^n = \bigoplus_{a + b = n} P_{-a}^b$ and
differential (as in the construction of the total complex associated
to a double complex) by
$$
\text{d}_P(x) = f_{-a}(x) + (-1)^a \text{d}_{P_{-a}}(x)
$$
for $x \in P_{-a}^b$. With these conventions $P$ is indeed a differential
graded $A$-module. Recalling that each $P_i$ has a two step filtration
$0 \to P_i' \to P_i \to P_i'' \to 0$ we set
$$
F_{2i}P = \bigoplus\nolimits_{i \geq j \geq 0} P_j
\subset
\bigoplus\nolimits_{i \geq 0} P_i = P
$$
and we add $P'_{i + 1}$ to $F_{2i}P$ to get $F_{2i + 1}$.
These are differential graded submodules and the successive quotients
are direct sums of shifts of $A$. By
Lemma \ref{lemma-target-graded-projective} we see that
the inclusions $F_iP \to F_{i + 1}P$ are admissible monomorphisms.
Finally, we have to show that the map $P \to M$ (given by the
augmentation $P_0 \to M$) is a quasi-isomorphism. This follows from
Homology, Lemma \ref{homology-lemma-good-resolution-gives-qis}.
\end{proof}

\begin{lemma}
\label{lemma-tensor-with-complex-homotopy}
The functor of Lemma \ref{lemma-tensor-with-complex} defines an exact functor
of triangulated categories
$K(\text{Mod}_{(E, \text{d})}) \to K(\mathcal{O})$.
\end{lemma}

\begin{proof}
The functor induces a functor between homotopy categories by
Lemma \ref{lemma-functorial}.
We have to show that $- \otimes_E K^\bullet$ transforms distinguished
triangles into distinguished triangles.
Suppose that $0 \to K \to L \to M \to 0$ is an admissible short
exact sequence of differential graded $E$-modules. Let $s : M \to L$ be
a graded $E$-module homomorphism which is left inverse to $L \to M$.
Then $s$ defines a map $M \otimes_E K^\bullet \to L \otimes_E K^\bullet$
of graded $\mathcal{O}$-modules (i.e., respecting $\mathcal{O}$-module
structure and grading, but not differentials)
which is left inverse to $L \otimes_E K^\bullet \to M \otimes_E K^\bullet$.
Thus we see that
$$
0 \to K \otimes_E K^\bullet \to L \otimes_E K^\bullet \to
M \otimes_E K^\bullet \to 0
$$
is a termwise split short exact sequences of complexes, i.e., a
defines a distinguished triangle in $K(\mathcal{O})$.
\end{proof}

\begin{lemma}
\label{lemma-existence-of-derived}
In the situation above. If the right derived functor $R\Hom(K^\bullet, -)$
of $\Hom(K^\bullet, -) : K(\mathcal{A}) \to D(\textit{Ab})$
is everywhere defined on $D(\mathcal{A})$, then we obtain a canonical exact
functor
$$
R\Hom(K^\bullet, -) : D(\mathcal{A}) \longrightarrow D(E, \text{d})
$$
of triangulated categories which reduces to the usual one  on taking
associated complexes of abelian groups.
\end{lemma}

\begin{proof}
Note that we have an associated functor
$K(\mathcal{A}) \to K(\text{Mod}_{(E, \text{d})})$ by
Lemma \ref{lemma-construction}.
We claim this functor is an exact functor of triangulated categories.
Namely, let $f : A^\bullet \to B^\bullet$ be a map of complexes
of $\mathcal{A}$. Then a computation shows that
$$
\Hom_{\text{Comp}^{dg}(\mathcal{A})}(K^\bullet, C(f)^\bullet)
=
C\left(
\Hom_{\text{Comp}^{dg}(\mathcal{A})}(K^\bullet, A^\bullet) \to
\Hom_{\text{Comp}^{dg}(\mathcal{A})}(K^\bullet, B^\bullet)
\right)
$$
where the right hand side is the cone in $\text{Mod}_{(E, \text{d})}$
defined earlier in this chapter.
This shows that our functor is compatible with cones, hence with
distinguished triangles. Let $X^\bullet$ be an object of $K(\mathcal{A})$.
Consider the category of quasi-isomorphisms $s : X^\bullet \to Y^\bullet$.
We are given that the functor
$(s : X^\bullet \to Y^\bullet) \mapsto \Hom_\mathcal{A}(K^\bullet, Y^\bullet)$
is essentially constant when viewed in $D(\textit{Ab})$.
But since the forgetful functor $D(E, \text{d}) \to D(\textit{Ab})$
is compatible with taking cohomology, the same thing is true in
$D(E, \text{d})$. This proves the lemma.
\end{proof}

\begin{lemma}
\label{lemma-property-I-K-injective}
Let $(A, \text{d})$ be a differential graded algebra.
Let $I$ be a differential graded $A$-module with property (I).
Then
$$
\Hom_{K(\text{Mod}_{(A, \text{d})})}(N, I) = 0
$$
for all acyclic differential graded $A$-modules $N$.
\end{lemma}

\begin{proof}
We will use that $K(\text{Mod}_{(A, \text{d})})$ is a triangulated
category (Proposition \ref{proposition-homotopy-category-triangulated}).
Let $F_\bullet$ be a filtration on $I$ as in property (I).
The short exact sequence of Lemma \ref{lemma-property-I-sequence}
produces a distinguished triangle. Hence by
Derived Categories, Lemma \ref{derived-lemma-representable-homological}
it suffices to show that
$$
\Hom_{K(\text{Mod}_{(A, \text{d})})}(N, I/F_iI) = 0
$$
for all acyclic differential graded $A$-modules $N$ and all $i$.
Each of the differential graded modules $I/F_iI$ has a finite filtration
by admissible monomorphisms, whose graded pieces are
products of $A^\vee[k]$. Thus it suffices to prove that
$$
\Hom_{K(\text{Mod}_{(A, \text{d})})}(N, A^\vee[k]) = 0
$$
for all acyclic differential graded $A$-modules $N$ and all $k$.
This follows from Lemma \ref{lemma-hom-into-shift-dual-free}
and the fact that $(-)^\vee$ is an exact functor.
\end{proof}

\begin{lemma}
\label{lemma-property-P-K-projective}
Let $(A, \text{d})$ be a differential graded algebra.
Let $P$ be a differential graded $A$-module with property (P).
Then
$$
\Hom_{K(\text{Mod}_{(A, \text{d})})}(P, N) = 0
$$
for all acyclic differential graded $A$-modules $N$.
\end{lemma}

\begin{proof}
We will use that $K(\text{Mod}_{(A, \text{d})})$ is a triangulated
category (Proposition \ref{proposition-homotopy-category-triangulated}).
Let $F_\bullet$ be a filtration on $P$ as in property (P).
The short exact sequence of Lemma \ref{lemma-property-P-sequence}
produces a distinguished triangle. Hence by
Derived Categories, Lemma \ref{derived-lemma-representable-homological}
it suffices to show that
$$
\Hom_{K(\text{Mod}_{(A, \text{d})})}(F_iP, N) = 0
$$
for all acyclic differential graded $A$-modules $N$ and all $i$.
Each of the differential graded modules $F_iP$ has a finite filtration
by admissible monomorphisms, whose graded pieces are direct sums
of shifts $A[k]$. Thus it suffices to prove that
$$
\Hom_{K(\text{Mod}_{(A, \text{d})})}(A[k], N) = 0
$$
for all acyclic differential graded $A$-modules $N$ and all $k$.
This follows from Lemma \ref{lemma-hom-from-shift-free}.
\end{proof}

\begin{lemma}
\label{lemma-right-resolution}
Let $(A, \text{d})$ be a differential graded algebra.
Let $M$ be a differential graded $A$-module. There exists a homomorphism
$M \to I$ of differential graded $A$-modules such that
\begin{enumerate}
\item $M \to I$ is a quasi-isomorphism, and
\item $I$ has property (I).
\end{enumerate}
\end{lemma}

\begin{proof}
Set $M = M_0$. We inductively choose short exact sequences
$$
0 \to M_i \to I_i \to M_{i + 1} \to 0
$$
where the maps $M_i \to I_i$ are chosen as in Lemma \ref{lemma-good-sub}.
This gives a ``resolution''
$$
0 \to M \to I_0 \xrightarrow{f_0} I_1 \xrightarrow{f_1} I_1 \to \ldots
$$
Denote $I$ the differential graded $A$-module with graded parts
$$
I^n = \prod\nolimits_{i \geq 0} I^{n - i}_i
$$
and differential defined by
$$
\text{d}_I(x) = f_i(x) + (-1)^i \text{d}_{I_i}(x)
$$
for $x \in I_i^{n - i}$. With these conventions $I$ is indeed a differential
graded $A$-module. Recalling that each $I_i$ has a two step filtration
$0 \to I_i' \to I_i \to I_i'' \to 0$ we set
$$
F_{2i}I^n = \prod\nolimits_{j \geq i} I^{n - j}_j
\subset
\prod\nolimits_{i \geq 0} I^{n - i}_i = I^n
$$
and we add a factor $I'_{i + 1}$ to $F_{2i}I$ to get $F_{2i + 1}I$.
These are differential graded submodules and the successive quotients
are products of shifts of $A^\vee$. By
Lemma \ref{lemma-source-graded-injective} we see that
the inclusions $F_{i + 1}I \to F_iI$ are admissible monomorphisms.
Finally, we have to show that the map $M \to I$ (given by the
augmentation $M \to I_0$) is a quasi-isomorphism. This follows from
Homology, Lemma \ref{homology-lemma-good-right-resolution-gives-qis}.
\end{proof}

\begin{lemma}
\label{lemma-good-quotient}
Let $(A, \text{d})$ be a differential graded algebra.
Let $M$ be a differential graded $A$-module. There exists a homomorphism
$P \to M$ of differential graded $A$-modules with the following
properties
\begin{enumerate}
\item $P \to M$ is surjective,
\item $\Ker(\text{d}_P) \to \Ker(\text{d}_M)$ is surjective, and
\item $P$ sits in an admissible short exact sequence
$0 \to P' \to P \to P'' \to 0$ where $P'$, $P''$ are direct sums
of shifts of $A$.
\end{enumerate}
\end{lemma}

\begin{proof}
Let $P_k$ be the free $A$-module with generators $x, y$ in degrees
$k$ and $k + 1$. Define the structure of a differential graded
$A$-module on $P_k$ by setting $\text{d}(x) = y$ and $\text{d}(y) = 0$.
For every element $m \in M^k$ there is a homomorphism
$P_k \to M$ sending $x$ to $m$ and $y$ to $\text{d}(m)$.
Thus we see that there is a surjection from a direct sum
of copies of $P_k$ to $M$. This clearly produces $P \to M$
having properties (1) and (3). To obtain property (2) note
that if $m \in \Ker(\text{d}_M)$ has degree $k$, then there is a map
$A[k] \to M$ mapping $1$ to $m$. Hence we can achieve (2) by adding
a direct sum of copies of shifts of $A$.
\end{proof}

\begin{lemma}
\label{lemma-target-graded-projective}
Let $(A, \text{d})$ be a differential graded algebra.
Let $M \to P$ be a surjective homomorphism of differential graded
$A$-modules. If $P$ is projective as a graded $A$-module, then
$M \to P$ is an admissible epimorphism.
\end{lemma}

\begin{proof}
This is immediate from the definitions.
\end{proof}

\begin{lemma}
\label{lemma-functorial}
Let $R$ be a ring. A functor $F : \mathcal{A} \to \mathcal{B}$
of differential graded categories over $R$ induces functors
$\text{Comp}(\mathcal{A}) \to \text{Comp}(\mathcal{B})$
and $K(\mathcal{A}) \to K(\mathcal{B})$.
\end{lemma}

\begin{proof}
Omitted.
\end{proof}

\begin{lemma}
\label{lemma-tensor-with-complex}
In the situation above there is a functor
$$
- \otimes_E K^\bullet :
\text{Mod}^{dg}_{(E, \text{d})}
\longrightarrow
\text{Comp}^{dg}(\mathcal{O})
$$
of differential graded categories. This functor sends $E$ to $K^\bullet$
and commutes with direct sums.
\end{lemma}

\begin{proof}
Let $M$ be a differential graded $E$-module. For every object $U$ of
$\mathcal{C}$ the complex $K^\bullet(U)$ is a left differential
graded $E$-module as well as a right $\mathcal{O}(U)$-module.
The actions commute, so we have a bimodule.
Thus, by the constructions in
Sections \ref{section-tensor-product} and \ref{section-bimodules}
we can form the tensor product
$$
M \otimes_E K^\bullet(U)
$$
which is a differential graded $\mathcal{O}(U)$-module, i.e., a complex
of $\mathcal{O}(U)$-modules. This construction is functorial with respect
to $U$, hence we can sheafify to get a complex of $\mathcal{O}$-modules
which we denote
$$
M \otimes_E K^\bullet
$$
Moreover, for each $U$ the construction determines a functor
$\text{Mod}^{dg}_{(E, \text{d})} \to \text{Comp}^{dg}(\mathcal{O}(U))$
of differential graded categories by Lemma \ref{lemma-tensor}.
It is therefore clear that we obtain a functor as stated in the lemma.
\end{proof}

\begin{lemma}
\label{lemma-construction}
Let $R$ be a ring. Let $\mathcal{A}$ be a differential graded category
over $R$. Let $x$ be an object of $\mathcal{A}$. Let
$$
(E, \text{d}) = \Hom_\mathcal{A}(x, x)
$$
be the differential graded $R$-algebra of endomorphisms of $x$.
We obtain a functor
$$
\mathcal{A} \longrightarrow \text{Mod}^{dg}_{(E, \text{d})},\quad
y \longmapsto \Hom_\mathcal{A}(x, y)
$$
of differential graded categories by letting $E$ act on
$\Hom_\mathcal{A}(x, y)$ via composition in $\mathcal{A}$.
This functor induces functors
$$
\text{Comp}(\mathcal{A}) \to \text{Mod}_{(A, \text{d})}
\quad\text{and}\quad
K(\mathcal{A}) \to K(\text{Mod}_{(A, \text{d})})
$$
by an application of Lemma \ref{lemma-functorial}.
\end{lemma}

\begin{proof}
This lemma proves itself.
\end{proof}

\begin{lemma}
\label{lemma-hom-into-shift-dual-free}
Let $(A, \text{d})$ be a differential graded algebra. Then we have
$$
\Hom_{\text{Mod}_{(A, \text{d})}}(M, A^\vee[k]) =
\Ker(\text{d} : (M^\vee)^k \to (M^\vee)^{k + 1})
$$
and
$$
\Hom_{K(\text{Mod}_{(A, \text{d})})}(M, A^\vee[k]) = H^k(M^\vee)
$$
as functors in the differential graded $A$-module $M$.
\end{lemma}

\begin{proof}
This is clear from the discussion above.
\end{proof}

\begin{lemma}
\label{lemma-property-I-sequence}
Let $(A, \text{d})$ be a differential graded algebra.
Let $I$ be a differential graded $A$-module. If $F_\bullet$
is a filtration as in property (I), then we obtain an
admissible short exact sequence
$$
0 \to I \to
\prod\nolimits I/F_iI \to
\prod\nolimits I/F_iI \to 0
$$
of differential graded $A$-modules.
\end{lemma}

\begin{proof}
Omitted. Hint: This is dual to Lemma \ref{lemma-property-P-sequence}.
\end{proof}

\begin{proposition}
\label{proposition-homotopy-category-triangulated}
Let $(A, \text{d})$ be a differential graded algebra. The homotopy category
$K(\text{Mod}_{(A, \text{d})})$ of differential graded $A$-modules with its
natural translation functors and distinguished triangles is a triangulated
category.
\end{proposition}

\begin{proof}
We know that $K(\text{Mod}_{(A, \text{d})})$ is a pre-triangulated category.
Hence it suffices to prove TR4 and to prove it we can use
Derived Categories, Lemma \ref{derived-lemma-easier-axiom-four}.
Let $K \to L$ and $L \to M$ be composable morphisms of
$K(\text{Mod}_{(A, \text{d})})$. By
Lemma \ref{lemma-sequence-maps-split} we may assume that
$K \to L$ and $L \to M$ are admissible monomorphisms.
In this case the result follows from
Lemma \ref{lemma-two-split-injections}.
\end{proof}

\begin{lemma}
\label{lemma-hom-from-shift-free}
Let $(A, d)$ be a differential graded algebra. Then we have
$$
\Hom_{\text{Mod}_{(A, \text{d})}}(A[k], M) =
\Ker(\text{d} : M^{-k} \to M^{-k + 1})
$$
and
$$
\Hom_{K(\text{Mod}_{(A, \text{d})})}(A[k], M) = H^{-k}(M)
$$
for any differential graded $A$-module $M$.
\end{lemma}

\begin{proof}
Immediate from the definitions.
\end{proof}

\begin{lemma}
\label{lemma-good-sub}
Let $(A, \text{d})$ be a differential graded algebra.
Let $M$ be a differential graded $A$-module. There exists a homomorphism
$M \to I$ of differential graded $A$-modules with the following
properties
\begin{enumerate}
\item $M \to I$ is injective,
\item $\Coker(\text{d}_M) \to \Coker(\text{d}_I)$ is injective,
and
\item $I$ sits in an admissible short exact sequence
$0 \to I' \to I \to I'' \to 0$ where $I'$, $I''$ are products
of shifts of $A^\vee$.
\end{enumerate}
\end{lemma}

\begin{proof}
We will use the functors $N \mapsto N^\vee$ (from left to right
differential graded modules and from right to left differential
graded modules) constructed in
Section \ref{section-modules-noncommutative-differential-graded}
and all of their properties.
For every $k \in \mathbf{Z}$ let $Q_k$ be the free left $A$-module with
generators $x, y$ in degrees $k$ and $k + 1$. Define the structure of a
left differential graded $A$-module on $Q_k$ by setting $\text{d}(x) = y$
and $\text{d}(y) = 0$. Arguing exactly as in the proof of
Lemma \ref{lemma-good-quotient} we find a surjection
$$
\bigoplus\nolimits_{i \in I} Q_{k_i} \longrightarrow M^\vee
$$
of left differential graded $A$-modules. Then we can consider the injection
$$
M \to (M^\vee)^\vee \to (\bigoplus\nolimits_{i \in I} Q_{k_i})^\vee =
\prod\nolimits_{i \in I} I_{k_i}
$$
where $I_k = Q_{-k}^\vee$ is the ``dual'' right differential graded $A$-module.
Further, the short exact sequence $0 \to A[-k - 1] \to Q_k \to A[-k] \to 0$
produces a short exact sequence
$0 \to A^\vee[k] \to I_k \to A^\vee[k + 1] \to 0$.

\medskip\noindent
The result of the previous paragraph produces $M \to I$
having properties (1) and (3). To obtain property (2), suppose
$\overline{m} \in \Coker(\text{d}_M)$ is a nonzero element of
degree $k$. Pick a map $\lambda : M^k \to \mathbf{Q}/\mathbf{Z}$
which vanishes on $\Im(M^{k - 1} \to M^k)$ but not on $m$. By
Lemma \ref{lemma-hom-into-shift-dual-free}
this corresponds to a homomorphism $M \to A^\vee[k]$ of
differential graded $A$-modules which does not vanish on $m$.
Hence we can achieve (2) by adding
a product of copies of shifts of $A^\vee$.
\end{proof}

\begin{lemma}
\label{lemma-source-graded-injective}
Let $(A, \text{d})$ be a differential graded algebra.
Let $I \to M$ be an injective homomorphism of differential graded
$A$-modules. If $I$ is graded injective, then
$I \to M$ is an admissible monomorphism.
\end{lemma}

\begin{proof}
This is immediate from the definitions.
\end{proof}

\begin{lemma}
\label{lemma-tensor}
Let $R$ be a ring. Let $(A, \text{d})$ and $(B, \text{d})$
be differential graded algebras over $R$. Let $N$ be a
differential graded $(A, B)$-bimodule. Then
$M \mapsto M \otimes_A N$ defines a functor
$$
- \otimes_A N :
\text{Mod}^{dg}_{(A, \text{d})}
\longrightarrow
\text{Mod}^{dg}_{(B, \text{d})}
$$
of differential graded categories. This functor induces functors
$$
\text{Mod}_{(A, \text{d})} \to \text{Mod}_{(B, \text{d})}
\quad\text{and}\quad
K(\text{Mod}_{(A, \text{d})}) \to K(\text{Mod}_{(B, \text{d})})
$$
by an application of Lemma \ref{lemma-functorial}.
\end{lemma}

\begin{proof}
Above we have seen how the construction defines a functor of underlying
graded categories. Thus it suffices to show that the construction is
compatible with differentials. Let $M$ and $M'$ be differential
graded $A$-modules and let $f : M \to M'$ be an $A$-module homomorphism
which is homogeneous of degree $n$. Then we have
$$
\text{d}(f) = \text{d}_{M'} \circ f - (-1)^n f \circ \text{d}_M
$$
On the other hand, we have
$$
\text{d}(f \otimes \text{id}_N) =
\text{d}_{M' \otimes_A N} \circ
(f \otimes \text{id}_N)
- (-1)^n
(f \otimes \text{id}_N) \circ \text{d}_{M \otimes_A N}
$$
Applying this to an element $x \otimes y$ with $x \in M$ and
$y \in N$ homogeneous we get
\begin{align*}
\text{d}(f \otimes \text{id}_N)(x \otimes y)
= &
\text{d}_{M'}(f(x)) \otimes y + (-1)^{n + \deg(x)}f(x) \otimes \text{d}_N(y) \\
& - (-1)^n f(\text{d}_M(x)) \otimes y
- (-1)^{n + \deg(x)}f(x) \otimes \text{d}_N(y) \\
= &
\text{d}(f) (x \otimes y)
\end{align*}
Thus we see that $\text{d}(f) \otimes \text{id}_N =
\text{d}(f \otimes \text{id}_N)$ and the proof is complete.
\end{proof}

\section{Bimodules}
\label{section-bimodules}

\section{Tensor product}
\label{section-tensor-product}

\begin{lemma}
\label{lemma-sequence-maps-split}
Let $(A, \text{d})$ be a differential graded algebra.
Let $L_1 \to L_2 \to \ldots \to L_n$
be a sequence of composable homomorphisms of
differential graded $A$-modules.
There exists a commutative diagram
$$
\xymatrix{
L_1 \ar[r] &
L_2 \ar[r] &
\ldots \ar[r] &
L_n \\
M_1 \ar[r] \ar[u] &
M_2 \ar[r] \ar[u] &
\ldots \ar[r] &
M_n \ar[u]
}
$$
in $\text{Mod}_{(A, \text{d})}$ such that each $M_i \to M_{i + 1}$
is an admissible monomorphism and each $M_i \to L_i$
is a homotopy equivalence.
\end{lemma}

\begin{proof}
The case $n = 1$ is without content.
Lemma \ref{lemma-make-injective} is the case $n = 2$.
Suppose we have constructed the diagram
except for $M_n$. Apply Lemma \ref{lemma-make-injective} to
the composition $M_{n - 1} \to L_{n - 1} \to L_n$.
The result is a factorization $M_{n - 1} \to M_n \to L_n$
as desired.
\end{proof}

\begin{lemma}
\label{lemma-two-split-injections}
Let $(A, \text{d})$ be a differential graded algebra. Suppose that
$\alpha : K \to L$ and $\beta : L \to M$ are admissible monomorphisms
of differential graded $A$-modules. Then there exist distinguished triangles
$(K, L, Q_1, \alpha, p_1, d_1)$, $(K, M, Q_2, \beta \circ \alpha, p_2, d_2)$
and $(L, M, Q_3, \beta, p_3, d_3)$ for which TR4 holds.
\end{lemma}

\begin{proof}
Say $\pi_1 : L \to K$ and $\pi_3 : M \to L$ are homomorphisms
of graded $A$-modules which are left inverse to $\alpha$ and $\beta$.
Then also $K \to M$ is an admissible monomorphism with left
inverse $\pi_2 = \pi_1 \circ \pi_3$. Let us write $Q_1$, $Q_2$
and $Q_3$ for the cokernels of $K \to L$, $K \to M$, and $L \to M$.
Then we obtain identifications (as graded $A$-modules)
$Q_1 = \Ker(\pi_1)$, $Q_3 = \Ker(\pi_3)$ and
$Q_2 = \Ker(\pi_2)$. Then $L = K \oplus Q_1$ and
$M = L \oplus Q_3$ as graded $A$-modules. This implies
$M = K \oplus Q_1 \oplus Q_3$. Note that $\pi_2 = \pi_1 \circ \pi_3$
is zero on both $Q_1$ and $Q_3$. Hence $Q_2 = Q_1 \oplus Q_3$.
Consider the commutative diagram
$$
\begin{matrix}
0 & \to & K & \to & L & \to & Q_1 & \to & 0 \\
  &     & \downarrow &     & \downarrow &     & \downarrow  & \\
0 & \to & K & \to & M & \to & Q_2 & \to & 0 \\
  &     & \downarrow &     & \downarrow &     & \downarrow  & \\
0 & \to & L & \to & M & \to & Q_3 & \to & 0
\end{matrix}
$$
The rows of this diagram are admissible short exact sequences, and
hence determine distinguished triangles by definition. Moreover
downward arrows in the diagram above are compatible with the chosen
splittings and hence define morphisms of triangles
$$
(K \to L \to Q_1 \to K[1])
\longrightarrow
(K \to M \to Q_2 \to K[1])
$$
and
$$
(K \to M \to Q_2 \to K[1])
\longrightarrow
(L \to M \to Q_3 \to L[1]).
$$
Note that the splittings $Q_3 \to M$ of the bottom sequence in the
diagram provides a splitting for the split sequence
$0 \to Q_1 \to Q_2 \to Q_3 \to 0$ upon composing with $M \to Q_2$.
It follows easily from this that the morphism $\delta : Q_3 \to Q_1[1]$
in the corresponding distinguished triangle
$$
(Q_1 \to Q_2 \to Q_3 \to Q_1[1])
$$
is equal to the composition $Q_3 \to L[1] \to Q_1[1]$.
Hence we get a structure as in the conclusion of axiom TR4.
\end{proof}

\section{Injective modules and differential graded algebras}
\label{section-modules-noncommutative-differential-graded}

\begin{lemma}
\label{lemma-make-injective}
Let $(A, \text{d})$ be a differential graded algebra.
Let $\alpha : K \to L$ be a homomorphism of differential graded
$A$-modules. There exists a factorization
$$
\xymatrix{
K \ar[r]^{\tilde \alpha} \ar@/_1pc/[rr]_\alpha &
\tilde L \ar[r]^\pi & L
}
$$
in $\text{Mod}_{(A, \text{d})}$ such that
\begin{enumerate}
\item $\tilde \alpha$ is an admissible monomorphism (see
Definition \ref{definition-admissible-ses}),
\item there is a morphism $s : L \to \tilde L$
such that $\pi \circ s = \text{id}_L$ and such that
$s \circ \pi$ is homotopic to $\text{id}_{\tilde L}$.
\end{enumerate}
\end{lemma}

\begin{proof}
The proof is identical to the proof of
Derived Categories, Lemma \ref{derived-lemma-make-injective}.
Namely, we set $\tilde L = L \oplus C(1_K)$ and we use elementary
properties of the cone construction.
\end{proof}

\begin{definition}
\label{definition-admissible-ses}
Let $(A, \text{d})$ be a differential graded algebra.
\begin{enumerate}
\item A homomorphism $K \to L$ of differential graded $A$-modules
is an {\it admissible monomorphism} if there exists a graded $A$-module
map $L \to K$ which is left inverse to $K \to L$.
\item A homomorphism $L \to M$ of differential graded $A$-modules
is an {\it admissible epimorphism} if there exists a graded $A$-module
map $M \to L$ which is right inverse to $L \to M$.
\item A short exact sequence $0 \to K \to L \to M \to 0$ of differential
graded $A$-modules is an {\it admissible short exact sequence}
if it is split as a sequence of graded $A$-modules.
\end{enumerate}
\end{definition}

\begin{lemma}
\label{lemma-fully-faithful-in-compact-case}
Let $(\mathcal{C}, \mathcal{O})$ be a ringed site.
Let $K^\bullet$ be a complex of $\mathcal{O}$-modules.
Assume
\begin{enumerate}
\item $K^\bullet$ represents a compact object of $D(\mathcal{O})$, and
\item $E = \Hom_{\text{Comp}^{dg}(\mathcal{O})}(K^\bullet, K^\bullet)$
computes the ext groups of $K^\bullet$ in $D(\mathcal{O})$.
\end{enumerate}
Then the functor
$$
- \otimes_E^\mathbf{L} K^\bullet :
D(E, \text{d})
\longrightarrow
D(\mathcal{O})
$$
of Lemma \ref{lemma-tensor-with-complex-derived} is fully faithful.
\end{lemma}

\begin{proof}
Because our functor has a left adjoint given by
$R\Hom(K^\bullet, -)$ by Lemma \ref{lemma-tensor-with-complex-hom-adjoint}
it suffices to show for a differential graded $E$-module $M$ that the map
$$
H^0(M) \longrightarrow
\Hom_{D(\mathcal{O})}(K^\bullet, M \otimes_E^\mathbf{L} K^\bullet)
$$
is an isomorphism. We may assume that $M = P$ is a differential graded
$E$-module which has property (P). Since $K^\bullet$ defines a
compact object, we reduce using
Lemma \ref{lemma-property-P-sequence}
to the case where $P$ has a finite filtration whose graded pieces
are direct sums of $E[k]$. Again using compactness we reduce
to the case $P = E[k]$. The assumption on $K^\bullet$ is that
the result holds for these.
\end{proof}

\end{document}
